\section{Arithmetic evaluation after averaging the shift}
\label{sec:averaged-energy}

We now square first and average the shift afterwards.  Let
\[
 G\in C_c^1((0,\infty)),\qquad G\geq0,
\]
with $I_0$ and $I_2$ as in \eqref{eq:moment-functionals}.
The shift sum below is over all integers; the support of $G(h/H)$ makes it
finite and positive for large $H$.

\subsection{The first-moment model}

\begin{proposition}[Long-shift first moment]
\label{prop:first-moment-model}
Fix $C\geq1$.  Uniformly for $X\leq H\leq X^C$, every real $t$, every
$D\geq1$, and every $A>0$,
\begin{equation}
 \frac1H\sum_hG(h/H)\cT_{h,D}(X;t)
 =I_0(G)\cM_D(X;t)+O_{A,C,W,G}\!\left(X(\log X)^{-A}\right).
 \label{eq:first-moment-model}
\end{equation}
Equivalently, the same normalized average of $\cR_{h,D}$ is
$O_A(X(\log X)^{-A})$.
\end{proposition}

\begin{proof}
For $r\in[\alpha X,\beta X]$, the smooth prime number theorem gives,
uniformly in $r$ and $H$ in the stated range,
\begin{equation}
 \sum_hG(h/H)(\Lambda(r+h)-1)
 \ll_{A,C,W,G}H(\log X)^{-A}.
 \label{eq:translated-smooth-PNT}
\end{equation}
Indeed, after the change of variable $q=r+h$, the weight
$G((q-r)/H)$ has scale $H$ and is supported on $q\asymp_{C,G}H$.
The classical zero-free-region prime number theorem is more than
sufficient.

The remaining source absolute mass satisfies
\[
 \sum_{\substack{n>D\\m\geq1}}|u(n)W(mn/X)|
 \ll_W X\sum_{n\ll_W X}\frac{1+\Lambda(n)}n+O_W(X)
 \ll_W X\log(2X).
\]
Insert \eqref{eq:translated-smooth-PNT}, request two additional powers of
$\log X$, and divide by $H$.  The Mellin phase has modulus one, so the
estimate is uniform in $t$.
\end{proof}

The coherent average in \cref{prop:first-moment-model} is not a theorem
about most shifts.  That requires the second moment developed next.

\subsection{The diagonal inputs}

\begin{lemma}[Translated target second moment]
\label{lem:target-second-moment}
Uniformly for $r\in[\alpha X,\beta X]$ and $X\leq H\leq X^C$,
\begin{equation}
 \sum_hG(h/H)(\Lambda(r+h)-1)^2
 =H I_0(G)\log H+O_{C,W,G}(H).
 \label{eq:target-second-moment}
\end{equation}
The same formula, with the same leading term, holds with
$(\Lambda(r+h)-1)^2$ replaced by $\Lambda(r+h)^2$.
\end{lemma}

\begin{proof}
Expand $(\Lambda-1)^2=\Lambda^2-2\Lambda+1$.  Smooth partial summation and
the prime number theorem give
\[
 \sum_hG(h/H)\Lambda(r+h)=H I_0(G)+O_A(H(\log H)^{-A})
\]
and
\[
 \sum_hG(h/H)=H I_0(G)+O_G(1).
\]
For the square, first restrict to prime target values.  One von Mangoldt
factor supplies the prime number theorem, while the other is
$\log H+O_{C,W,G}(1)$ on the support.  Thus the prime contribution is
$H I_0(G)\log H+O_{C,W,G}(H)$.  Proper prime powers contribute
$O_{C,W,G}(H^{1/2}(\log H)^2)$, which is absorbed.  Combining the three
pieces proves the centered formula.
\end{proof}

\begin{lemma}[Source diagonal]
\label{lem:source-diagonal}
Uniformly for $1\leq D\leq X^{1-\eps}$,
\begin{equation}
 \begin{aligned}
 \cS_D(X)
 &:={}
 \sum_{\substack{n>D\\m\geq1}}
 u(n)^2|W(mn/X)|^2\\
 &=\frac12X I_2(W)
 \bigl((\log X)^2-(\log D)^2\bigr)
 +O_W(X\log(2X)).
 \end{aligned}
 \label{eq:source-diagonal}
\end{equation}
\end{lemma}

\begin{proof}
Euler summation gives, uniformly in $n$,
\begin{equation}
 \sum_{m\geq1}|W(mn/X)|^2
 =\frac Xn I_2(W)+O_W(1).
 \label{eq:W-square-Euler}
\end{equation}
The elementary weighted second moment
\begin{equation}
 \sum_{n\leq y}\frac{(\Lambda(n)-1)^2}{n}
 =\frac12(\log y)^2-\log y+O(1)
 \label{eq:source-weighted-second-moment}
\end{equation}
follows by expanding the square, using the classical zero-free-region
prime number theorem for
$\sum\Lambda(n)^2/n$ and $\sum\Lambda(n)/n$, and summing proper prime
powers absolutely.  Also
$\sum_{n\leq y}(\Lambda(n)-1)^2\ll y\log(2y)$.  Insert
\eqref{eq:W-square-Euler}, subtract the ranges up to $D$, and absorb the
fixed support endpoints into $O_W(X\log X)$.
\end{proof}

\subsection{The same-ray off-diagonal ledger}

\begin{lemma}[Two-target upper sieve]
\label{lem:two-target-upper-sieve}
Let $r\ne r'$ lie in $[\alpha X,\beta X]$.  Uniformly for
$X\leq H\leq X^C$,
\begin{equation}
 \begin{aligned}
 &\sum_hG(h/H)
 \left|(\Lambda(r+h)-1)(\Lambda(r'+h)-1)\right|\\
 &\qquad\ll_{C,W,G}
 H\prod_{\substack{p\mid r-r'\\p>2}}\frac{p-1}{p-2}
 \ll_{C,W,G}H\log\log(3X).
 \end{aligned}
 \label{eq:two-target-upper-sieve}
\end{equation}
The same bound holds if either centered target factor is replaced by
$\Lambda$.
\end{lemma}

\begin{proof}
The constant and one-prime terms cost $O_G(H)$.  If both target values are
prime, the Selberg upper-bound sieve for the two linear forms $h+r$ and
$h+r'$ gives
\[
 \#\{h\asymp_G H:h+r,\ h+r'\ \hbox{prime}\}
 \ll_G\frac{H}{(\log H)^2}
 \prod_{\substack{p\mid r-r'\\p>2}}\frac{p-1}{p-2}.
\]
Restoring the two logarithmic weights proves the first bound for the
prime-prime part.  If either target is a proper prime power, the aggregate
contribution is $O_G(H^{1/2}(\log H)^2)$.

Finally,
\[
 \prod_{\substack{p\mid q\\p>2}}\frac{p-1}{p-2}
 \ll\log\log(3|q|)
\]
by the standard maximal-order comparison with the product over the
smallest primes; here $0<|r-r'|\ll_W X$.  This proves
\eqref{eq:two-target-upper-sieve}.  See, for example,
\citet[Chapter~6]{HalberstamRichert1974} for the two-linear-forms upper
sieve.
\end{proof}

\begin{lemma}[Same-ray source off-diagonal]
\label{lem:same-ray-source-ledger}
For every $D\geq1$,
\begin{equation}
 \begin{aligned}
 \cL_D(X):={}&
 \sum_{\substack{(a,b)=1}}
 \sum_{\substack{k\ne\ell\\ak,a\ell>D}}
 |u(ak)u(a\ell)|\\
 &\quad\times
 |W(abk^2/X)W(ab\ell^2/X)|
 \ll_W X(\log(2X))^2.
 \end{aligned}
 \label{eq:same-ray-source-ledger}
\end{equation}
\end{lemma}

\begin{proof}
For a fixed ray use the set $K_{a,b}$ from
\cref{prop:fixed-ray-energy-bound} and put
$x_k=|u(ak)W(abk^2/X)|$.  Then
\[
 \sum_{k\ne\ell}x_kx_\ell
 \leq |K_{a,b}|\sum_kx_k^2
 \ll_W\sum_{k\in K_{a,b}}k x_k^2
\]
by \eqref{eq:ray-cardinality}.  Summing over $b$ and applying
\eqref{eq:number-of-b-rays} yields
\[
 \cL_D(X)
 \ll_W X\sum_{ak^2\ll_W X}\frac{u(ak)^2}{ak}.
\]
Now use \cref{lem:source-ray-square-ledger}.
\end{proof}

\subsection{The averaged ray-energy theorem}

\begin{theorem}[Averaged-shift ray energy]
\label{thm:averaged-ray-energy}
Fix $\eps>0$ and $C\geq1$.  Uniformly for
\[
 1\leq D\leq X^{1-\eps},
 \qquad X\leq H\leq X^C,
\]
one has
\begin{equation}
 \begin{aligned}
 \sum_hG(h/H)\cE_{h,D}(X)
 ={}&\frac12HX I_0(G)I_2(W)\log H\\
 &\times\bigl((\log X)^2-(\log D)^2\bigr)\\
 &+O_{\eps,C,W,G}\!\left(
 HX(\log(2X))^2\log\log(3X)\right).
 \end{aligned}
 \label{eq:averaged-ray-energy}
\end{equation}
The same asymptotic holds for $\cE^{\mathrm{raw}}_{h,D}(X)$.
\end{theorem}

\begin{proof}
Expand the square in \eqref{eq:ray-energy} and interchange the finite
sums.  For $k=\ell$, the target argument is $r=abk^2$ and
\cref{lem:target-second-moment} gives
\[
 H I_0(G)\log H\sum_{(a,b)=1}\sum_k
 u(ak)^2|W(abk^2/X)|^2+O_G(H\cS_D(X)).
\]
The double source sum is exactly $\cS_D(X)$: the substitution
$n=ak,m=bk$ is unique.  Inserting \cref{lem:source-diagonal} produces the
main term in \eqref{eq:averaged-ray-energy}.  All diagonal errors are
$O_{\eps,C,W,G}(HX(\log X)^2)$.

For $k\ne\ell$, the two target arguments
\[
 r=abk^2,
 \qquad r'=ab\ell^2
\]
are distinct.  Apply \cref{lem:two-target-upper-sieve}, take absolute
values of the source coefficients, and then use
\cref{lem:same-ray-source-ledger}.  The total off-diagonal contribution is
\[
 \ll_{C,W,G}H\log\log(3X)\,\cL_D(X)
 \ll_{C,W,G}HX(\log X)^2\log\log(3X).
\]
This proves the centered theorem.  For the raw energy, the target
diagonal $\Lambda^2$ has the same leading term by
\cref{lem:target-second-moment}, and the same upper-sieve ledger handles
the off-diagonal.
\end{proof}

\begin{remark}[Where the saving comes from]
No cancellation is asserted in an individual same-ray prime correlation.
The target diagonal contributes $H\log H$, while the absolute
off-diagonal upper sieve contributes only $H\log\log X$.  The source
diagonal is of order $X\log^2X$, and the entire source off-diagonal ledger
has the same order.  The extra logarithm on the target diagonal is the
sole reason an asymptotic results.
\end{remark}
